\section{Introduction}
\label{sec:introduction}

Urban Air Mobility (UAM) is emerging as a promising extension of future transportation systems. To enable the safe integration of large numbers of electric vertical take-off and landing (eVTOL) aircraft into low-altitude airspace, structured traffic-management frameworks such as the European U-space~\cite{eu_commission_2017,easa2020uspace} and the U.S. UAS Traffic Management (UTM) concept~\cite{kopardekar2016utm} have been proposed. These initiatives emphasize service-oriented architectures and trajectory-based operations to support scalable and safe airspace integration~\cite{sesar_conops_2019,sesar_master_plan_2025}. It should be noted that the U-space framework itself does not prescribe corridors: U-space services are provided within designated regions or volumes, and unmanned aircraft system (UAS) operators can plan their trajectories within those volumes. Nevertheless, UAM corridors are envisioned by some operators for early UAM implementation, with U-space services provided for a larger volume covering all UAM corridors of a region, as investigated, for example, in the SESAR project EUREKA~\cite{sesar_eureka_project}. This paper therefore addresses structured corridor operations within U-space rather than treating U-space as a corridor concept.

A central requirement in such environments is computationally efficient path planning. In dense urban scenarios, trajectory requests must be processed while respecting separation requirements and structured corridor constraints. Search-based approaches such as A* have been widely investigated for UAS routing~\cite{sahawneh2016_3d_path_planning,a_star_2d3d_comparison_utm}, alongside 4D grid-based conflict management strategies~\cite{acevedo_4d_grids_u_space}. Deconfliction in multi-dimensional trajectory spaces has also been studied in trajectory-based operations frameworks~\cite{NDMap,enea_porretta_4d_tbo_comparison}, and surveys summarize the broader range of UAS path-planning methods~\cite{ahmad2023_uav_path_planning_survey}.

In operational U-space systems, path planning is not an isolated routing problem. A selected route must be converted into a time-parameterized 4D trajectory, checked against operational constraints, and integrated into a service-oriented traffic-management workflow. Monolithic implementations of these functions may limit extensibility and independent scaling. Microservice architectures offer a modular alternative by separating concerns such as topology provision, trajectory generation, deconfliction, and authorization into dedicated services~\cite{dragoni2017microservices,matlekovic2020uav}. However, quantitative evaluations of service-oriented U-space planning systems remain limited, particularly with respect to latency decomposition, computational scaling with path distance, and the evaluated behavior of integrated conflict-checking and re-planning loops.

This paper presents a service-oriented 4D path-planning architecture for structured corridor operations within U-space and evaluates both planner-stage computational scaling and conflict-enabled behavior. Two existing DLR tools are used as building blocks. The Copter Flight Management System (CFMS) is a modular experimental flight-management and 4D trajectory generation system for unmanned multicopter and other VTOL aircraft that converts a phase-structured waypoint sequence into a dense, kinematically consistent, time-parameterized trajectory~\cite{morscheck2021cfms}. NDMap is a high-performance conflict-detection library for multi-dimensional objects that indexes trajectories in a combined space--time state space and reports separation violations between them~\cite{NDMap}. The architecture separates graph-based route planning, waypoint synthesis, 4D trajectory generation, NDMap-based deconfliction, and U-space Service Provider (USSP) interaction into distinct service responsibilities. The planner stage combines A*-based routing on a vertiport corridor graph with CFMS-based trajectory generation~\cite{morscheck2021cfms}. A CesiumJS-based frontend enables interactive visualization of the generated time-dependent trajectories~\cite{cesiumjs_software}.

The evaluated planning pipeline follows four main steps. First, A* selects a shortest corridor route on the vertiport graph. Second, waypoint synthesis augments this route with departure, climb, cruise, descent, and landing structure. Third, CFMS converts the structured waypoint sequence into a dense time-parameterized 4D trajectory. Fourth, when conflict checking is enabled, NDMap checks the candidate trajectory against previously accepted trajectories and conflict intervals are fed back to the Planning Service for request-local lateral re-routing. Each planning request in this paper corresponds to a single operator request for trajectory generation and flight-plan approval for one origin--destination vertiport pair.

The contributions of this work are:
\begin{itemize}
    \item A service-oriented U-space path-planning architecture that separates topology provision, planning, trajectory generation, deconfliction, and USSP interaction responsibilities.
    \item A CFMS-compatible waypoint-synthesis pipeline that converts A*-selected corridor routes into phase-structured inputs for 4D trajectory generation.
    \item A planner-stage latency decomposition into A* path computation, waypoint synthesis, and CFMS invocation with trajectory generation.
    \item An NDMap-based conflict-checking loop with request-local lateral re-routing by graph-edge removal and A* re-planning.
    \item A sequential Hamburg vertiport corridor-network evaluation of both planner-stage computational scaling and conflict-resolution-enabled behavior.
\end{itemize}

The remainder of the paper is organized as follows. Section~\ref{sec:related} reviews related work, Section~\ref{sec:methodology} describes the system design, Section~\ref{sec:experimental_setup} outlines the evaluation setup, Section~\ref{sec:results} presents the results, Section~\ref{sec:discussion} discusses the findings and their limitations, and Section~\ref{sec:conclusion} concludes the paper.
